\section{Task Formulation}
\label{sec:problem}

An agent accumulates an interaction history $H=(s_1,\ldots,s_T)$ of dated
dialogue sessions. A memory construction process maps this history to a
memory $\mem$ comprising the source passages and representations derived
from them. Given a question $q$, retrieval constructs a context $\ctx$
under a budget $b$, from which a reader produces an answer:
\begin{equation}
 \ctx=\mathrm{Retrieve}(q,\mem;b),\qquad
 \hat a=\mathrm{Read}(q,\ctx).
 \label{eq:task}
\end{equation}
The budget bounds the evidence delivered to the reader, so the goal of
retrieval is to select evidence that \emph{supports} a correct answer,
rather than records that merely resemble the question.

We consider questions whose relational requirements can be expressed as
a conjunctive query (\cref{sec:preliminaries}). Under this view,
retrieval amounts to finding witnesses of the query in memory, with the
temporal scope and expected form of the answer shaping which witnesses
are preferred. The formulation thus distinguishes two responsibilities:
\emph{interpreting} the question as a query, and \emph{satisfying} that
query over stored facts. Witnesses make the second responsibility
explicit and checkable; their adequacy for the original question still
depends on the interpretation and on the stored facts, a point we return
to in \cref{sec:discussion}.
